%-------------------------
% Lebenslauf in LaTeX (DE)
% Based off of: https://github.com/sb2nov/resume
%-------------------------

\documentclass[a4paper,11pt]{article}

\usepackage[T1]{fontenc}
\usepackage[utf8]{inputenc}
\usepackage{latexsym}
\usepackage[empty]{fullpage}
\usepackage{titlesec}
\usepackage{marvosym}
\usepackage[usenames,dvipsnames]{color}
\usepackage{verbatim}
\usepackage{enumitem}
\usepackage[hidelinks]{hyperref}
\usepackage{fancyhdr}
\usepackage[ngerman]{babel}
\usepackage{tabularx}
\usepackage{needspace}
\input{glyphtounicode}

\pagestyle{fancy}
\fancyhf{}
\fancyfoot{}
\renewcommand{\headrulewidth}{0pt}
\renewcommand{\footrulewidth}{0pt}

% Seitenränder
\addtolength{\oddsidemargin}{-0.5in}
\addtolength{\evensidemargin}{-0.5in}
\addtolength{\textwidth}{1in}
\addtolength{\topmargin}{-.5in}
\addtolength{\textheight}{1.0in}

\urlstyle{same}
\raggedbottom
\raggedright
\setlength{\tabcolsep}{0in}

% Formatierung der Abschnitte
\titleformat{\section}{
  \vspace{-4pt}\scshape\raggedright\large
}{}{0em}{}[\color{black}\titlerule \vspace{-5pt}]

% Maschinenlesbares / ATS-freundliches PDF
\pdfgentounicode=1

%-------------------------
% Benutzerdefinierte Befehle
\newcommand{\resumeItem}[1]{
  \item\small{
    {#1 \vspace{-2pt}}
  }
}

\newcommand{\resumeSubheading}[4]{
  \Needspace{5\baselineskip}
  \vspace{-2pt}\item
  \textbf{#1}\hfill #2\\
  \textit{\small #3}\hfill \textit{\small #4}\\
  \vspace{-7pt}
}

\newcommand{\resumeProjectHeading}[2]{
    \item
    \begin{tabular*}{0.97\textwidth}{l@{\extracolsep{\fill}}r}
      \small#1 & #2 \\
    \end{tabular*}\vspace{-7pt}
}

\newcommand{\resumeOneLineHeading}[2]{
  \item
  \small #1\hfill #2\\[-1pt]
}

\renewcommand\labelitemii{$\vcenter{\hbox{\tiny$\bullet$}}$}

\newcommand{\resumeSubHeadingListStart}{\begin{itemize}[leftmargin=0.15in, label={}]}
\newcommand{\resumeSubHeadingListEnd}{\end{itemize}}
\newcommand{\resumeItemListStart}{\begin{itemize}}
\newcommand{\resumeItemListEnd}{\end{itemize}\vspace{-5pt}}

%-------------------------------------------
%%%%%%  LEBENSLAUF  %%%%%%%%%%%%%%%%%%%%%%%%

\begin{document}

\begin{center}
    \textbf{\Large \scshape Stanislav Smirnov} \\ \vspace{1pt}
    \small Köln, Deutschland $|$ +49 176 4345 8829 $|$ \href{mailto:mr.stassmirnoff@hotmail.com}{\underline{mr.stassmirnoff@hotmail.com}} $|$
    \href{https://www.linkedin.com/in/smirnov-stanislav/}{\underline{linkedin.com/in/smirnov-stanislav}} $|$
    \href{https://github.com/mrStasSmirnoff}{\underline{github.com/mrStasSmirnoff}}
\end{center}

%-----------KURZPROFIL-----------
\section{Kurzprofil}
Senior Data Scientist mit Schwerpunkt Payments und Fraud sowie Erfahrung in der Leitung von 0-to-1-ML- und Optimierungsprodukten in den Bereichen Routing, Retries, Decline Recovery, Tokenisierung und Fraud Detection. Ich verbinde ML und Analytics mit KPI-Verantwortung, Stakeholder-Alignment und produktiver Umsetzung, um Autorisierungsraten zu verbessern, Umsatz zurückzugewinnen und Betrug zu reduzieren.

%-----------ERFOLGE-----------
\section{Erfolge}
\begin{itemize}[leftmargin=*, noitemsep, topsep=0pt, partopsep=0pt]
    \item Steigerte die Erfolgsquote im Fallback-Routing um den Faktor \textbf{4} und den zurückgewonnenen Umsatz um den Faktor \textbf{3} durch Analyse der PSP-Performance im nordamerikanischen Traffic und Neugestaltung eines schwach performenden Rerouting-Pfads.
    \item Konzipierte und skalierte die ML-Lösung \textbf{Dynamic Retries} von der Idee bis zum Produktivbetrieb und verbesserte die Retry-Erfolgsrate für die relevanten Transaktionen durch KPI-basierte Experimente und Rollout um \textbf{25--50\%}.
    \item Baute ein produktives Fraud-ML-Modell neu auf und steigerte den Fraud Recall von \textbf{${\sim}$44\% auf ${\sim}$72\%} bei unveränderter \textbf{False-Positive-Rate von 1\%}; der PR-AUC stieg von \textbf{0.51 auf 0.80}.
    \item Steuerte die cross-funktionale Umsetzung mit Product, Engineering und Fraud Prevention von Priorisierung und KPI-Definition bis zu Rollout und Evaluation.
\end{itemize}
\vspace{-2pt}

%-----------KOMPETENZEN-----------
\section{Kompetenzen}
 \begin{itemize}[leftmargin=0.15in, label={}]
    \small{\item{
     \textbf{Führung \& Umsetzung}{: Stakeholder-Management, cross-funktionale Führung, End-to-End-Verantwortung, strukturierte Problemdefinition, Mentoring von Junior Data Scientists} \\
     \textbf{Produkt \& Strategie}{: KPI-Definition, Priorisierung, Entscheidungsvorlagen / PRDs, Rollout-Planung, Launch-Evaluation, Roadmap-/Backlog-Definition, Experimentdesign} \\
     \textbf{Payments \& Fraud}{: Autorisierungs-/Decline-Analyse, PSP-Performance, Routing, Retries, Tokenisierung (Network Tokens), Payment-Optimierung, Fraud Detection} \\
     \textbf{Programmierung \& ML/AI}{: Python, SQL, PySpark, klassisches ML, Deep Learning (CNNs, LSTMs), LLMs} \\
     \textbf{Cloud \& Data}{: AWS, Azure, Data Vault 2.0, ETL-Entwicklung} \\
     \textbf{Tools}{: Git, Docker, Airflow, FastAPI, SageMaker, Databricks, TensorFlow, Keras}
    }}
 \end{itemize}

%-----------BERUFSERFAHRUNG-----------
\section{Berufserfahrung}
  \resumeSubHeadingListStart

    \resumeSubheading
      {Senior Data Scientist, Payments}{April 2024 -- heute}
      {Cleverbridge AG}{Köln, DE}
      \resumeItemListStart
        \resumeItem{Leite Initiativen zur Payment-Optimierung und Fraud-ML in den Bereichen Routing, Retries, Autorisierung, Tokenisierung und Fraud Detection. Ich verantworte den Weg von Opportunity Discovery und Entscheidungsvorlage über Stakeholder-Alignment bis zum Rollout, definiere KPIs und bewerte die Ergebnisse nach dem Launch.}

        \resumeItem{Steigerte die Autorisierungsrate im Fallback-Rerouting um den Faktor \textbf{4} (von 2--3\% auf 11\%), indem ich einen schwachen Routing-Pfad im nordamerikanischen Traffic identifizierte und den Wechsel zu einer leistungsfähigeren Alternative anstieß; der zurückgewonnene Umsatz in diesem Traffic-Segment stieg um den Faktor \textbf{3}.}

        \resumeItem{Initiierte und leitete den Neuaufbau eines produktiven Fraud-ML-Modells, das zwei Jahre lang nicht nachtrainiert worden war. Gemeinsam mit Fraud Prevention überarbeitete ich die Fraud-Labeling-Logik, trainierte und deployte das Modell und steigerte den Offline-Fraud-Recall von \textbf{${\sim}$44\% auf ${\sim}$72\%} bei unveränderter \textbf{False-Positive-Rate von 1\%} (PR-AUC: \textbf{0.51 auf 0.80}). Die Evaluation verglich das bestehende Modell mit dem Score eines externen Anbieters und unterstützte die Build-vs-Buy-Entscheidung.}

        \resumeItem{Verantworte die Weiterentwicklung von \textbf{Dynamic Retries} und nutze Analysen auf BIN-, Regions-, PSP- und Network-Token-Ebene, um zusätzliche Routing- und Retry-Use-Cases zu priorisieren.}

        \resumeItem{Führe und mentore Junior Data Scientists in Problemdefinition, Modellentwicklung, Experimentdesign und produktionsnaher Analyse; verbessere die Lieferqualität durch Code-Reviews und strukturiertes Feedback.}
      \resumeItemListEnd

    \resumeSubheading
      {Data Scientist}{Dezember 2021 -- April 2024}
      {Cleverbridge AG}{Köln, DE}
      \resumeItemListStart
        \resumeItem{Konzipierte und brachte \textbf{Dynamic Retries}, ein ML-basiertes Produkt zur Optimierung des Retry-Timings, in Produktion. Verfasste den Produktvorschlag, stimmte Engineering, Product und Business aufeinander ab und leitete die Umsetzung von PoC über MVP bis zum Produktiv-Rollout.}

        \resumeItem{Steigerte die Erfolgsraten für die relevanten Retries um \textbf{25--50\%} und verantwortete KPI-Definition, Messdesign sowie die Weiterentwicklung auf Basis der Ergebnisse nach dem Rollout.}

        \resumeItem{Entwickelte eine hybride Retry-Strategie aus regelbasiertem Scheduling nach Land, Währung, Monatstag und Uhrzeit sowie ML-basierter Timing-Optimierung zur Auswahl von Retry-Zeitfenstern mit hoher Conversion.}

        \resumeItem{Implementierte und betrieb die Lösung auf AWS mit Python-Microservices, SageMaker und MWAA, einschließlich automatisierter Trainings- und Inference-Workflows für einen zuverlässigen Produktivbetrieb.}

        \resumeItem{Entwarf das ML-Framework hinter \textbf{CleverAutomations} und entwickelte Churn-Reduction-Modelle sowie automatisierte Workflows, die abwanderungsgefährdete Kunden identifizierten und gezielte Kampagnen auslösten.}
      \resumeItemListEnd

    \Needspace{12\baselineskip}
    \resumeSubheading
      {Data Engineer}{September 2019 -- Dezember 2021}
      {Cleverbridge AG}{Köln, DE}
      \resumeItemListStart
        \resumeItem{Leitete Design und Implementierung eines skalierbaren Data Warehouse auf Azure nach Data Vault 2.0 und ermöglichte damit zuverlässiges Reporting und Self-Service Analytics im Unternehmen.}

        \resumeItem{Entwickelte und optimierte ETL-Pipelines mit Azure Data Factory, Databricks und T-SQL zur Verarbeitung von Millionen Datensätzen pro Tag aus mehreren Quellsystemen.}

        \resumeItem{Steuerte Schemaänderungen und Migrationen mit Flyway und stellte konsistente, versionierte Deployments über mehrere Umgebungen sicher.}

        \resumeItem{Entwickelte Python-Microservices und Airflow-Workflows für Data Products und containerisierte sie mit Docker für einen robusten und wartbaren Betrieb.}
      \resumeItemListEnd

    \resumeSubheading
      {Machine Learning Intern}{Mai 2017 -- Juli 2017}
      {KOSTAL Group}{Dortmund, DE}
      \resumeItemListStart
        \resumeItem{Trainierte ein neuronales Netz zur Gestenerkennung mit Keras und TensorFlow und erreichte \textbf{80\%} Validierungsgenauigkeit für die fünf wichtigsten Gesten auf Basis statischer Bilddaten und eines CNN.}

        \resumeItem{Verbesserte die Robustheit des Modells durch Datenvorverarbeitung und Augmentation.}
      \resumeItemListEnd

    \resumeSubheading
      {Netzoptimierungsingenieur}{Mai 2012 -- März 2015}
      {MegaFon}{Republik Baschkortostan, Russland}
      \resumeItemListStart
        \resumeItem{Leitete ein \textbf{Netzoptimierungsteam} aus drei Ingenieuren.}

        \resumeItem{Reduzierte die CS Call Drop Rate im 2G-Netz um \textbf{0.2\%}.}

        \resumeItem{Steigerte die Inter-RAT Handover Success Rate im UMTS-Netz um \textbf{0.8\%} und verbesserte damit die Gesamtperformance des Netzes.}
      \resumeItemListEnd

  \resumeSubHeadingListEnd

%-----------PROJEKTE-----------
\section{Projekte}
  \resumeSubHeadingListStart

    \resumeProjectHeading
      {\textbf{Objekterkennung in Kunstgemälden (Masterarbeit)} $|$ \emph{Python, Deep Learning, CNN}}{2017 -- 2018}
      \resumeItemListStart
        \resumeItem{Entwickelte einen eigenen CNN-basierten Ansatz zur Objekterkennung in Kunstgemälden und verbesserte die Detektionsgenauigkeit gegenüber bestehenden Methoden um \textbf{18\%}.}

        \resumeItem{Baute den Trainings- und Evaluationsworkflow auf und führte eine systematische Fehleranalyse zur Weiterentwicklung des Modells durch.}
      \resumeItemListEnd

  \resumeSubHeadingListEnd

%-----------AUSBILDUNG-----------
\section{Ausbildung}
  \resumeSubHeadingListStart

    \resumeSubheading
      {Universität Paderborn}{Paderborn, DE}
      {Master of Science in Electrical Systems Engineering (Schwerpunkt Signal Processing)}{2015 -- 2018}
      \resumeItemListStart
        \resumeItem{\textbf{Masterarbeit:} Enhancing Object Detection in Fine Art Paintings Using Deep Learning Techniques (\textbf{+18\%} Detektionsgenauigkeit gegenüber der Baseline).}

        \resumeItem{\textbf{Wissenschaftliche Hilfskraft:} statistische Bildverarbeitung; Modellierung plasmonischer Nanopartikel in COMSOL; Programmierung von Lie-Gruppen in GAP.}
      \resumeItemListEnd

  \resumeSubHeadingListEnd

%-----------PUBLIKATIONEN & ZERTIFIZIERUNGEN-----------
\section{Publikationen \& Zertifizierungen}
  \begin{itemize}[leftmargin=0.15in, label={}, itemsep=2pt, topsep=0pt]

    \resumeOneLineHeading
      {\textbf{\href{https://ieeexplore.ieee.org/abstract/document/9089828}{\underline{Deep Learning for Object Detection in Fine Art Paintings}}} $|$ \emph{IEEE Xplore}}{2020}

    \resumeOneLineHeading
      {\textbf{\href{https://www.sciencedirect.com/science/article/abs/pii/S1359645420308491}{\underline{Nonlinear Dielectric Properties of Random Paraelectric-Dielectric Composites}}} $|$ \emph{Acta Materialia}}{2021}

    \resumeOneLineHeading
      {\textbf{\href{https://grow.cleverbridge.com/blog/failed-payment-recovery-dynamic-retries}{\underline{Dynamic Retries: failed payment recovery}}} $|$ \emph{Fachartikel / Technical Writing}}{}

    \resumeOneLineHeading
      {\textbf{AWS Certified Cloud Practitioner}}{}

  \end{itemize}

%-----------SPRACHEN-----------
\section{Sprachen}
 \begin{itemize}[leftmargin=0.15in, label={}]
    \small{\item{
      Russisch (Muttersprache) $|$ Englisch (fließend) $|$ Deutsch (B2)
    }}
 \end{itemize}

%-----------WEITERE INFORMATIONEN-----------
\section{Weitere Informationen}
 \begin{itemize}[leftmargin=0.15in, label={}]
    \small{\item{
      \textbf{Ehrenamt}: Regionalverband Baschkortostan der Russischen Union Junger Wissenschaftler. \\
      \textbf{Auszeichnung}: Dankschreiben des CEO der PJSC ``MegaFon'' in der Republik Baschkortostan für herausragende Leistungen.
    }}
 \end{itemize}

\end{document}